\documentclass[10pt,letterpaper]{article}
% ---------------------------------------------------------
% PAQUETES
% ---------------------------------------------------------
\usepackage[utf8]{inputenc}
\usepackage[T1]{fontenc}
\usepackage[spanish]{babel}
\usepackage{lmodern}
\usepackage[
    top=0.55in,
    bottom=0.55in,
    left=0.70in,
    right=0.70in
]{geometry}

\usepackage{xcolor}
\usepackage{enumitem}
\usepackage{titlesec}
\usepackage{hyperref}

% ---------------------------------------------------------
% COLORES
% ---------------------------------------------------------
\definecolor{NavyBlue}{HTML}{17365D}
\definecolor{DarkGray}{HTML}{333333}
\definecolor{MediumGray}{HTML}{666666}

% ---------------------------------------------------------
% CONFIGURACIÓN GENERAL
% ---------------------------------------------------------
\pagestyle{empty}

\setlength{\parindent}{0pt}
\setlength{\parskip}{0pt}

\hypersetup{
    colorlinks=true,
    urlcolor=NavyBlue,
    linkcolor=NavyBlue
}

% ---------------------------------------------------------
% SECCIONES
% ---------------------------------------------------------
\titleformat{\section}
    {\large\bfseries\color{NavyBlue}}
    {}
    {0pt}
    {}
    [\vspace{-5pt}\color{NavyBlue}\rule{\textwidth}{0.6pt}]

\titlespacing*{\section}
    {0pt}
    {9pt}
    {5pt}

% ---------------------------------------------------------
% LISTAS
% ---------------------------------------------------------
\setlist[itemize]{
    leftmargin=1.35em,
    itemsep=1.5pt,
    topsep=2pt,
    parsep=0pt,
    partopsep=0pt
}

% ---------------------------------------------------------
% COMANDOS PERSONALIZADOS
% ---------------------------------------------------------

\newcommand{\jobtitle}[1]{
    {\bfseries #1}
}

\newcommand{\company}[1]{
    {\bfseries\color{DarkGray} #1}
}

\newcommand{\jobdate}[1]{
    {\color{MediumGray} #1}
}

% ---------------------------------------------------------
% DOCUMENTO
% ---------------------------------------------------------
\begin{document}

% =========================================================
% ENCABEZADO
% =========================================================

\begin{center}

    {\fontsize{24}{28}\selectfont
    \bfseries\color{NavyBlue}
    ANA KAREN CANALES SANTOS}

    \vspace{3pt}

    {\large
    \bfseries
    Desarrolladora Java \textbar{} Ingeniera de Soporte de Aplicaciones
    \textbar{} L2/L3}

    \vspace{5pt}

    \small
Tulancingo de Bravo, Hidalgo, México
\quad | \quad
\href{mailto:ana.karen.canales.santos@gmail.com}{ana.karen.canales.santos@gmail.com}
\quad | \quad
\href{https://www.linkedin.com/in/anakarencanalessantos/}{LinkedIn}
\quad | \quad
\href{https://github.com/iHannahBeZombie}{GitHub}

\end{center}

% =========================================================
% PERFIL PROFESIONAL
% =========================================================

\section{Perfil Profesional}

Desarrolladora Java e Ingeniera de Soporte de Aplicaciones con 9 años de experiencia en desarrollo, soporte L2/L3 y operación de aplicaciones críticas en ambientes productivos. Especializada en el diagnóstico y resolución de incidencias en aplicaciones Java y APIs, análisis de logs, identificación de causa raíz (RCA) y coordinación con equipos técnicos y de negocio para resolver problemas de alto impacto y garantizar la continuidad operativa de los servicios.

Sólidos conocimientos en Java, Spring, REST APIs, SQL y Linux/Unix, así como en monitoreo y herramientas de observabilidad. Enfocada en mantener la estabilidad y disponibilidad de los servicios, automatizar tareas operativas e implementar soluciones que reduzcan la recurrencia de incidentes y mejoren la operación de las aplicaciones.

% ============================================
% HABILIDADES TÉCNICAS
% ============================================
\section{Habilidades Técnicas}

\begin{itemize}

    \item \textbf{Lenguajes:}
    Java, SQL, JavaScript, Bash, HTML, CSS, XML, JSON.

    \item \textbf{Frameworks y librerías Java:}
    Spring, Spring Boot, Spring MVC, Spring Data JPA, Hibernate,
    Spring Security, JavaServer Faces (JSF), Lombok, Jackson.

    \item \textbf{Frontend:}
    Angular, React, Bootstrap, jQuery.

    \item \textbf{APIs y servicios web:}
    REST APIs, SOAP / Web Services, Swagger, OpenAPI,
    Postman, SoapUI, Insomnia, cURL, OAuth 2.0, JWT.

    \item \textbf{Bases de datos:}
    Oracle, MySQL, SQL Server, H2, MongoDB;
    Oracle SQL Developer, MySQL Workbench, DataGrip, Navicat.

    \item \textbf{Cloud y almacenamiento:}
    Google Cloud Platform (GCP), Google Cloud Storage (GCS), BigQuery.

    \item \textbf{Control de versiones y construcción:}
    Git, GitHub, GitLab, Bitbucket, SVN, Maven, npm, Yarn.

    \item \textbf{Testing y calidad:}
    JUnit, Mockito, AssertJ, Spring Boot Test, SonarQube.

    \item \textbf{CI/CD y DevOps:}
    Jenkins, Docker, Kubernetes, pipelines CI/CD.

    \item \textbf{Monitoreo y observabilidad:}
    Splunk, Grafana, OpenSearch, Prometheus, xMatters.

    \item \textbf{Sistemas operativos:}
    Windows, Linux, Unix.

    \item \textbf{IDEs y herramientas de desarrollo:}
    Eclipse, IntelliJ IDEA, Spring Tool Suite, NetBeans,
    Visual Studio Code, Visual Studio, Notepad++, Sublime Text, Vim.

    \item \textbf{Documentación y diagramación:}
    Confluence, Draw.io, Lucidchart, Microsoft Visio, UML;
    documentación técnica, diagramas de arquitectura, runbooks,
    guías de procedimiento y manuales de usuario y administración.

    \item \textbf{Gestión de proyectos y servicios:}
    Jira, ServiceNow, Trello;
    Incident Management, Problem Management, Change Management,
    Root Cause Analysis (RCA), troubleshooting y soporte L2/L3.

    \item \textbf{Arquitectura y prácticas de desarrollo:}
    Arquitectura monolítica, microservicios, arquitectura hexagonal,
    SOLID, Design Patterns, Dependency Injection, Constructor Injection,
    Clean Code, Code Review, Agile y Scrum.

    \item \textbf{Integración y procesamiento de datos:}
    Procesamiento de archivos, CSV, Excel, FTP/FTPS,
    integración entre sistemas, APIs y servicios web.

\end{itemize}

% =========================================================
% EXPERIENCIA PROFESIONAL
% =========================================================

\section{Experiencia Profesional}

% ---------------------------------------------------------
% WALMART JULIO 2024-MAYO 2025
% ---------------------------------------------------------


\jobtitle{Ingeniera de Software II} \textbar{}
\company{Walmart}
\hfill \textbar{}
\jobdate{Julio 2024 -- Mayo 2025}

\begin{itemize}

    \item Proporcioné soporte de producción L2--L3 para aplicaciones Java y APIs críticas para el negocio en ambientes Linux de alta disponibilidad, asegurando la estabilidad, disponibilidad y cumplimiento de los compromisos de SLA.

    \item Desarrollé y mantuve componentes de aplicaciones utilizando Java, aplicando conocimientos de programación para diagnosticar problemas en producción, analizar el comportamiento de las aplicaciones, identificar defectos e implementar soluciones permanentes.

    \item Diagnostiqué y resolví incidentes complejos de producción mediante análisis de logs, consultas SQL, troubleshooting a nivel de código, análisis de causa raíz (RCA) y colaboración con equipos de ingeniería y soporte.

    \item Realicé consultas SQL para investigar incidentes de producción, validar datos de las aplicaciones, identificar inconsistencias y contribuir al análisis de causa raíz y resolución de incidentes.

    \item Monitoreé la salud de las aplicaciones, el rendimiento de los sistemas y los logs mediante Grafana, OpenSearch y Splunk, realizando revisiones periódicas de salud e identificando de manera proactiva anomalías y posibles incidentes.

    \item Gestioné el ciclo completo de incidentes, incluyendo triage, priorización, investigación, escalamiento, comunicación con stakeholders, resolución y análisis posterior al incidente, particularmente en eventos de alta severidad P1/P2.

    \item Coordiné con equipos multidisciplinarios y distribuidos geográficamente para impulsar la resolución de incidentes, asegurando una comunicación efectiva, ownership claro y la restauración oportuna de los servicios dentro de los compromisos de SLA.

    \item Realicé troubleshooting en ambientes Linux mediante herramientas básicas de línea de comandos, inspección de logs, monitoreo de procesos y verificaciones de salud del sistema para investigar problemas relacionados con las aplicaciones y la infraestructura.

    \item Trabajé con MongoDB para la investigación y troubleshooting de datos, incluyendo consultas y validación de información de las aplicaciones como apoyo al análisis de incidentes de producción.

    \item Adquirí experiencia práctica con Docker y Kubernetes, brindando soporte a aplicaciones Java contenerizadas y realizando troubleshooting básico de contenedores, deployments, pods y configuraciones de las aplicaciones.

    \item Implementé correcciones de código Java, hotfixes, cambios de configuración y actividades de remediación en producción, siguiendo procesos de gestión de cambios, gobierno de releases y despliegues con mínima interrupción del servicio.

    \item Analicé incidentes recurrentes y defectos de las aplicaciones para identificar oportunidades de mejora en código, configuración, datos y procesos operativos, contribuyendo a reducir la recurrencia de incidentes y mejorar la estabilidad de las plataformas.

    \item Automaticé actividades operativas y de soporte repetitivas mediante scripting, reduciendo el esfuerzo manual, mejorando los tiempos de respuesta e incrementando la eficiencia operativa.

    \item Creé y mantuve documentación técnica, incluyendo runbooks, guías de troubleshooting, procedimientos de reprocesamiento y artículos de la base de conocimiento en Confluence, mejorando la autonomía del equipo y la consistencia en la resolución de incidentes.

\end{itemize}

% ---------------------------------------------------------
% AGILETHOUGHT 2022-2023
% ---------------------------------------------------------

\jobtitle{Desarrolladora de Soporte Java} \textbar{}
\company{AgileThought -- Citibanamex}
\hfill \textbar{}
\jobdate{2022 -- 2023}

\begin{itemize}

\item Proporcioné soporte de aplicaciones y producción L2--L3 para sistemas backend desarrollados con Java y Spring Boot, realizando troubleshooting de APIs RESTful y módulos backend para garantizar la alta disponibilidad, estabilidad y operación confiable de los servicios en producción.

\item Desarrollé y mantuve componentes backend utilizando Java y Spring Boot, aplicando conocimientos de programación y debugging para investigar el comportamiento de las aplicaciones, identificar defectos e implementar soluciones permanentes en producción.

\item Investigué y resolví incidentes de producción mediante análisis de logs, debugging de aplicaciones, consultas SQL y análisis de causa raíz (RCA), colaborando con equipos de desarrollo y soporte para restaurar los servicios y prevenir la recurrencia de incidentes.

\item Brindé soporte al ciclo de vida de las aplicaciones desde una perspectiva operativa, incluyendo soporte a despliegues, validación en producción, resolución de incidentes, estabilización de releases y monitoreo posterior a los despliegues.

\item Monitoreé la salud de las aplicaciones y la infraestructura en ambientes AWS mediante CloudWatch, analizando logs, métricas y alertas para identificar de manera proactiva problemas de rendimiento, fallas de las aplicaciones y posibles incidentes.

\item Realicé investigaciones sobre bases de datos mediante consultas SQL para validar el comportamiento de las aplicaciones, analizar datos de producción, identificar inconsistencias y apoyar las actividades de troubleshooting y análisis de causa raíz.

\item Brindé soporte a despliegues y cambios en producción mediante pipelines de CI/CD, realizando validaciones de releases, troubleshooting de problemas durante los despliegues y soporte a actividades de rollback, minimizando la interrupción de las operaciones del negocio.

\item Realicé troubleshooting en ambientes Linux mediante operaciones de línea de comandos, inspección de logs, monitoreo de procesos y verificaciones de salud del sistema para investigar problemas relacionados con aplicaciones y despliegues.

\item Colaboré con equipos multidisciplinarios en ambientes Agile/Scrum para gestionar incidentes, coordinar releases, comunicar hallazgos técnicos y dar seguimiento a los problemas hasta su resolución dentro de los tiempos establecidos y compromisos de SLA.

\item Demostré ownership sobre los incidentes de producción, dando seguimiento de principio a fin a las actividades de investigación, resolución, validación, documentación y seguimiento posterior al incidente.

\item Creé y mantuve documentación operativa, incluyendo runbooks, guías de troubleshooting, procedimientos de despliegue y flujos de operación de los sistemas, mejorando la eficiencia del soporte, la transferencia de conocimiento y la autonomía del equipo.

\end{itemize}

% ---------------------------------------------------------
% INE Marzo 2021 - Noviembre 2021
% ---------------------------------------------------------

\jobtitle{Desarrolladora Java} \textbar{}
\company{INE}
\hfill \textbar{}
\jobdate{Marzo 2021 -- Noviembre 2021}

\begin{itemize}

\item Analicé requerimientos y reglas de operación, participando en el diseño y desarrollo de módulos orientados a optimizar la funcionalidad y el desempeño de las aplicaciones.

\item Desarrollé aplicaciones backend utilizando Java y Spring Boot, implementando componentes con Spring MVC para la exposición de servicios RESTful.

\item Implementé servicios RESTful para operaciones CRUD (Create, Read, Update, Delete), desarrollando los flujos necesarios para consultar, registrar, actualizar y eliminar información.

\item Utilicé Spring Data JPA e Hibernate para implementar la persistencia de datos y realizar consultas sobre bases de datos relacionales, asegurando la consistencia de la información.

\item Participé en el desarrollo de un servicio de gestión de usuarios con mecanismos de autenticación y autorización basados en roles, fortaleciendo la seguridad de las operaciones de la aplicación.

\item Implementé y configuré mecanismos de logging mediante Log4j2 para facilitar el seguimiento de eventos, diagnóstico de errores y troubleshooting de la aplicación.

\item Trabajé con las funcionalidades principales de Spring Boot, incluyendo auto-configuración, gestión de dependencias, logging y configuración de la aplicación mediante YAML, facilitando el desarrollo y despliegue de los servicios.

\end{itemize}

% ---------------------------------------------------------
% ACCENTURE Marzo 2020 - Noviembre 2020
% ---------------------------------------------------------

\jobtitle{Desarrolladora Java} \textbar{}
\company{Accenture -- Banorte}
\hfill \textbar{}
\jobdate{Marzo 2020 -- Noviembre 2020}

\begin{itemize}

\item Realicé un inventario integral de servicios y participé en la migración de transferencias de archivos desde FTP hacia FTPS, fortaleciendo la seguridad de la información y el cumplimiento de estándares de la organización.

\item Diseñé e implementé un plan de configuración para actualizar parámetros de servidores y habilitar protocolos de transferencia de archivos seguros y confiables.

\item Ejecuté actividades de validación y pruebas sobre los procesos de transferencia de archivos, verificando el manejo correcto de la información, la integridad de los datos y la continuidad operativa durante la migración.

\item Utilicé SQL Server y SVN para apoyar las actividades técnicas y de mantenimiento de los componentes involucrados en la transición, contribuyendo a mantener la disponibilidad de los servicios y minimizar riesgos operativos.

\item Trabajé con componentes desarrollados en .NET (C# 2010) como parte de las actividades de soporte y transición de los sistemas involucrados.

\item Colaboré con equipos de infraestructura y soporte para coordinar la implementación, validación e integración de los cambios de protocolo en un ambiente bancario crítico para el negocio.

\end{itemize}

% ---------------------------------------------------------
% IRONBIT DIGITAL EXPERIENCES
% ---------------------------------------------------------
\jobtitle{Desarrolladora Java} \textbar{}
\company{Ironbit Digital Experiences}
\hfill \textbar{}
\jobdate{Noviembre 2019 -- Febrero 2020}

\begin{itemize}

\item Desarrollé soluciones para la comparación automática de documentos, mejorando la capacidad de procesamiento y automatizando actividades relacionadas con la validación de información.

\item Diseñé y estructuré bases de datos, definiendo tablas y campos para organizar y almacenar la información requerida por las aplicaciones.

\item Participé en el análisis y rediseño de requerimientos para servicios REST, adaptando su funcionalidad a las necesidades del proyecto y mejorando los flujos de intercambio de información.

\item Desarrollé e implementé consultas utilizando JPA para facilitar la interacción entre las aplicaciones Java y la base de datos Oracle.

\item Desarrollé servicios REST para la consulta y almacenamiento de información, mejorando el intercambio de datos entre los diferentes componentes del sistema.

\item Implementé integraciones mediante servicios REST externos para consultar información de otras aplicaciones, contribuyendo a la integración efectiva entre sistemas.

\item Participé en la planeación y distribución de actividades técnicas dentro del equipo, contribuyendo a optimizar la carga de trabajo y el cumplimiento de los objetivos del proyecto.

\item Utilicé Java, JPA, Oracle y SVN, así como técnicas de modelado y diseño de bases de datos, para desarrollar e implementar las soluciones requeridas por el proyecto.

\end{itemize}

% ---------------------------------------------------------
% ACI Group
% ---------------------------------------------------------
\jobtitle{Senior Specialist A -- Desarrolladora Java} \textbar{}
\company{ACI Group}
\hfill \textbar{}
\jobdate{Mayo 2019 -- Octubre 2019}

\begin{itemize}

\item Lideré el desarrollo de servicios para la plataforma web del proyecto ``SGP Process Management System'' de Telcel, orientado a la gestión y mantenimiento de sitios de infraestructura.

\item Participé en el análisis de requerimientos, asegurando la alineación de las soluciones técnicas con las necesidades funcionales del proyecto y de los stakeholders.

\item Desarrollé nuevos módulos y funcionalidades para el sistema de generación de reportes, mejorando la capacidad de consulta y presentación de información.

\item Desarrollé e implementé consultas SQL y Stored Procedures para facilitar la manipulación, consulta y extracción de información desde la base de datos Oracle 11.

\item Participé en la corrección de bugs y en la implementación de mejoras funcionales y técnicas, contribuyendo a mantener la operación eficiente y estable del sistema.

\item Brindé atención y soporte técnico al cliente, resolviendo dudas, recopilando requerimientos y facilitando la comunicación entre los stakeholders y el equipo de desarrollo.

\item Participé activamente en la planeación y ceremonias de la metodología SAFe, incluyendo PI Planning, Backlog Refinement y System Demos, dentro de un ambiente de integración continua.

\item Desarrollé soluciones utilizando Groovy y Grails Framework, GSP, HTML y CSS, gestionando el código fuente mediante SVN.

\end{itemize}


% ---------------------------------------------------------
% QUARKSOFT
% ---------------------------------------------------------

\jobtitle{Desarrolladora Java} \textbar{}
\company{Quarksoft}
\hfill \textbar{}
\jobdate{Diciembre 2017 -- Mayo 2019}

\begin{itemize}

\item Participé en el desarrollo de la plataforma web ``DIS Platform of the Social Investment Directorate'' para Nacional Monte de Piedad, enfocada en la gestión de procesos de inversión social.

\item Desarrollé servicios REST para la gestión de información y funcionalidades relacionadas con la inversión social, realizando previamente el análisis detallado de requerimientos para asegurar su alineación con las necesidades del proyecto.

\item Desarrollé e implementé consultas nativas, HQL y operaciones mediante JPA/Hibernate para facilitar el acceso, consulta y manipulación de información en la base de datos MySQL.

\item Desarrollé y mantuve componentes backend utilizando Java y Spring Boot, gestionando las dependencias y el proceso de construcción del proyecto mediante Maven.

\item Colaboré en un equipo multidisciplinario junto con diferentes empresas de consultoría dentro de un ambiente de integración continua, coordinando actividades técnicas y dando seguimiento a los requerimientos del proyecto.

\item Brindé soporte al cliente, resolviendo dudas, recopilando necesidades y facilitando la comunicación entre los diferentes equipos y stakeholders involucrados en el proyecto.

\item Participé en la introducción y capacitación sobre el uso de módulos de la aplicación, facilitando la adopción y utilización de las funcionalidades desarrolladas.

\item Realicé pruebas funcionales de caja negra (black box testing) para validar el correcto funcionamiento de los servicios desarrollados y contribuir a la calidad de las soluciones entregadas.

\end{itemize}

% =========================================================
% CONECTA
% =========================================================
\jobtitle{Desarrolladora Backend Java} \textbar{}
\company{Conecta}
\hfill \textbar{}
\jobdate{Marzo 2017 -- Diciembre 2017}

\begin{itemize}

\item Desarrollé la aplicación web ``API-REST Web Application Monitor'', orientada al monitoreo y generación de reportes sobre la concurrencia de aplicaciones, incluyendo el seguimiento de solicitudes y el registro de resultados exitosos o fallidos.

\item Diseñé y desarrollé servicios REST, definiendo URIs para facilitar el acceso a la información y documentando la API para mejorar su comprensión, consumo y mantenimiento.

\item Implementé mecanismos de control de acceso mediante la definición de roles, permitiendo la consulta y modificación de información de acuerdo con la jerarquía organizacional.

\item Implementé Hazelcast como plataforma distribuida de datos, contribuyendo al rendimiento y escalabilidad de la aplicación.

\item Participé en el levantamiento y análisis de requerimientos, asegurando que las funcionalidades desarrolladas respondieran a las necesidades del cliente.

\item Elaboré diagramas UML, incluyendo diagramas de casos de uso, clases y secuencia, para modelar y documentar la estructura, comportamiento y flujos principales del sistema.

\item Desarrollé pruebas unitarias para validar el correcto funcionamiento de los componentes y servicios, contribuyendo a la calidad y confiabilidad de la aplicación.

\item Trabajé con Java, Spring, Hibernate y MySQL para el desarrollo de los componentes backend y la interacción con la base de datos.

\end{itemize}

% =========================================================
% BST MÉXICO
% =========================================================
\jobtitle{Desarrolladora Backend Java} \textbar{}
\company{BST México}
\hfill \textbar{}
\jobdate{Octubre 2015 -- Febrero 2017}

\begin{itemize}

\item Participé en el desarrollo de servicios y funcionalidades para la plataforma web ``My Habitat'', orientada a la gestión, administración y control contable de activos inmobiliarios.

\item Desarrollé y mantuve funcionalidades de la plataforma, implementando nuevas características, realizando mejoras y corrigiendo bugs para optimizar el funcionamiento y desempeño del sistema.

\item Participé en actividades de pruebas y análisis de experiencia de usuario, identificando oportunidades de mejora y validando el correcto funcionamiento de las funcionalidades desarrolladas.

\item Realicé pruebas unitarias y documenté sus resultados como parte de las actividades de aseguramiento de la calidad del software.

\item Administré y mantuve el código fuente del proyecto, contribuyendo a una gestión ordenada y eficiente durante el ciclo de desarrollo.

\item Elaboré documentación técnica y manuales de uso para facilitar la comprensión, operación y adopción de la plataforma por parte del personal y los usuarios.

\item Participé en actividades de capacitación y presentación del sistema, proporcionando orientación sobre el uso de las funcionalidades y contribuyendo a su adopción por parte de los usuarios.

\item Colaboré en la atención de necesidades funcionales y técnicas de la plataforma, participando en el análisis, desarrollo, pruebas y mejora continua de las soluciones implementadas.

\end{itemize}
% =========================================================
% EDUCACIÓN
% =========================================================

\section{Educación}

\textbf{Ingeniería en Sistemas Computacionales} \textbar{}
Universidad Politécnica de Tulancingo
\hfill
\jobdate{2012 -- 2015}

% =========================================================
% IDIOMAS
% =========================================================

\section{Idiomas}

\textbf{Español:} Nativo

\textbf{Inglés:} B2

\end{document}